% GENERATED FILE. Edit canonical lessons, then run npm run generate:books.
% canonical-source-hash: fnv1a64:03e1a9daf03dfa68
% canonical-lessons: SA-C489-vyasanam, SA-C489-samyamah, SA-C489-dhumapanam, SA-C489-parityajati, SA-C489-madyam

\chapter{Habits and Self-Control}
\label{ch:sa-habits-and-self-control}
\hlchaptermodality{489}

\begin{chapteropening}
\textbf{By the end of this chapter:} \emph{I can talk about a bad habit and giving it up.}

Complete the last lesson of chapter 489: I can talk about a bad habit and giving it up.
\end{chapteropening}

\section[\texorpdfstring{\sk{व्यसनम्} (vyasanam)}{व्यसनम् (vyasanam)}]{\sk{व्यसनम्} (vyasanam) — a vice; an addiction}

\label{lesson:SA-C489-vyasanam}



\begin{quote}
Before the new one: say the Sanskrit for sleepless, then the Sanskrit for a three-hour watch.
\end{quote}

\subsection*{You'll want to know: \sk{व्यसनम्}}
\textbf{\sk{व्यसनम्}} — \emph{vyasanam} — \textquotedblleft{}a vice; an addiction\textquotedblright{}.

\begin{quote}
\sk{द्यूतम्} \sk{एकं} \sk{व्यसनम्} \sk{अस्ति।}
\end{quote}

\begin{quote}
\emph{dyūtam ekaṃ vyasanam asti.}
\end{quote}

\begin{quote}
Gambling is a vice.
\end{quote}

A neuter noun.

\begin{cousinweb}
\textbf{\sk{व्यसनम्}} (\emph{vyasanam}) is built from \emph{vi} and the root √as, \textquotedblleft{}to throw\textquotedblright{}: throwing oneself into a thing; MW gives \textquotedblleft{}attachment or addiction to\textquotedblright{} and \textquotedblleft{}evil passion, vice\textquotedblright{} (MW).
\end{cousinweb}

\begin{grammarlens}[title={What you've built}]
One more word for describing what you lived and felt.
\end{grammarlens}

\subsection*{Your turn}
\begin{itemize}
\raggedright
  \item \emph{Say it:} \emph{vyasanam}
  \item \emph{Say it:} \emph{vyasanam}, once more
  \item \emph{Say it:} \emph{dyūtam ekaṃ vyasanam asti.} — Gambling is a vice.
  \item \emph{Recall:} say the Sanskrit for sleepless, then the Sanskrit for a three-hour watch, then say \emph{vyasanam} again
\end{itemize}

\begin{tcolorbox}[breakable,colback=teal!4,colframe=teal!35!black,title={Before you move on}]
What does \sk{व्यसनम्} mean? (\textquotedblleft{}A vice; an addiction\textquotedblright{}.) Say it once more.
\end{tcolorbox}
\section[\texorpdfstring{\sk{संयमः} (saṃyamaḥ)}{संयमः (saṃyamaḥ)}]{\sk{संयमः} (saṃyamaḥ) — self-control}

\label{lesson:SA-C489-samyamah}



\begin{quote}
Before the new one: say the Sanskrit for a three-hour watch, then the Sanskrit for a vice.
\end{quote}

\subsection*{You'll want to know: \sk{संयमः}}
\textbf{\sk{संयमः}} — \emph{saṃyamaḥ} — \textquotedblleft{}self-control\textquotedblright{}.

\begin{quote}
\sk{संयमेन} \sk{व्यसनं} \sk{त्यजतु।}
\end{quote}

\begin{quote}
\emph{saṃyamena vyasanaṃ tyajatu.}
\end{quote}

\begin{quote}
Give up the vice through self-control.
\end{quote}

A masculine noun.

\begin{cousinweb}
\textbf{\sk{संयमः}} (\emph{saṃyamaḥ}) is built from \emph{sam} and the root √yam, \textquotedblleft{}to hold\textquotedblright{}; MW gives \textquotedblleft{}restraint, control of the senses, self-control\textquotedblright{} (MW).
\end{cousinweb}

\begin{grammarlens}[title={What you've built}]
One more word for describing what you lived and felt.
\end{grammarlens}

\subsection*{Your turn}
\begin{itemize}
\raggedright
  \item \emph{Say it:} \emph{saṃyamaḥ}
  \item \emph{Say it:} \emph{saṃyamaḥ}, once more
  \item \emph{Say it:} \emph{saṃyamena vyasanaṃ tyajatu.} — Give up the vice through self-control.
  \item \emph{Recall:} say the Sanskrit for a three-hour watch, then the Sanskrit for a vice, then say \emph{saṃyamaḥ} again
\end{itemize}

\begin{tcolorbox}[breakable,colback=teal!4,colframe=teal!35!black,title={Before you move on}]
What does \sk{संयमः} mean? (\textquotedblleft{}Self-control\textquotedblright{}.) Say it once more.
\end{tcolorbox}
\section[\texorpdfstring{\sk{धूमपानम्} (dhūmapānam)}{धूमपानम् (dhūmapānam)}]{\sk{धूमपानम्} (dhūmapānam) — smoking}

\label{lesson:SA-C489-dhumapanam}



\begin{quote}
Before the new one: say the Sanskrit for a vice, then the Sanskrit for self-control.
\end{quote}

\subsection*{You'll want to know: \sk{धूमपानम्}}
\textbf{\sk{धूमपानम्}} — \emph{dhūmapānam} — \textquotedblleft{}smoking\textquotedblright{}.

\begin{quote}
\sk{धूमपानं} \sk{फुप्फुसं} \sk{पीडयति।}
\end{quote}

\begin{quote}
\emph{dhūmapānaṃ phupphusaṃ pīḍayati.}
\end{quote}

\begin{quote}
Smoking harms the lungs.
\end{quote}

A neuter noun.

\begin{cousinweb}
\textbf{\sk{धूमपानम्}} (\emph{dhūmapānam}) is built from \textbf{\sk{धूमः}}, \textquotedblleft{}smoke\textquotedblright{}, and \emph{pāna}, \textquotedblleft{}drinking\textquotedblright{}, from the root √pā of \textbf{\sk{पिबति}}; MW gives \textquotedblleft{}inhalation of smoke or vapour\textquotedblright{} and \textquotedblleft{}tobacco-smoking\textquotedblright{} (MW).
\end{cousinweb}

\begin{grammarlens}[title={What you've built}]
One more word for describing what you lived and felt.
\end{grammarlens}

\subsection*{Your turn}
\begin{itemize}
\raggedright
  \item \emph{Say it:} \emph{dhūmapānam}
  \item \emph{Say it:} \emph{dhūmapānam}, once more
  \item \emph{Say it:} \emph{dhūmapānaṃ phupphusaṃ pīḍayati.} — Smoking harms the lungs.
  \item \emph{Recall:} say the Sanskrit for a vice, then the Sanskrit for self-control, then say \emph{dhūmapānam} again
\end{itemize}

\begin{tcolorbox}[breakable,colback=teal!4,colframe=teal!35!black,title={Before you move on}]
What does \sk{धूमपानम्} mean? (\textquotedblleft{}Smoking\textquotedblright{}.) Say it once more.
\end{tcolorbox}
\section[\texorpdfstring{\sk{परित्यजति} (parityajati)}{परित्यजति (parityajati)}]{\sk{परित्यजति} (parityajati) — gives up entirely}

\label{lesson:SA-C489-parityajati}



\begin{quote}
Before the new one: say the Sanskrit for self-control, then the Sanskrit for smoking.
\end{quote}

\subsection*{You'll want to know: \sk{परित्यजति}}
\textbf{\sk{परित्यजति}} — \emph{parityajati} — \textquotedblleft{}gives up entirely\textquotedblright{}.

\begin{quote}
\sk{पिता} \sk{धूमपानं} \sk{परित्यजति।}
\end{quote}

\begin{quote}
\emph{pitā dhūmapānaṃ parityajati.}
\end{quote}

\begin{quote}
Father gives up smoking.
\end{quote}

Past (he …): \sk{सः} \sk{धूमपानं} \sk{परित्यक्तवान्}, \emph{saḥ dhūmapānaṃ parityaktavān}.

\begin{cousinweb}
\textbf{\sk{परित्यजति}} (\emph{parityajati}) is built from \emph{pari}, \textquotedblleft{}around, wholly\textquotedblright{}, and the root √tyaj, \textquotedblleft{}to leave\textquotedblright{}, the root of \textbf{\sk{त्यजति}}; MW gives \textquotedblleft{}to leave, quit, abandon, give up\textquotedblright{} (MW).
\end{cousinweb}

\begin{grammarlens}[title={What you've built}]
One more word for describing what you lived and felt.
\end{grammarlens}

\subsection*{Your turn}
\begin{itemize}
\raggedright
  \item \emph{Say it:} \emph{parityajati}
  \item \emph{Say it:} \emph{parityajati}, once more
  \item \emph{Say it:} \emph{pitā dhūmapānaṃ parityajati.} — Father gives up smoking.
  \item \emph{Recall:} say the Sanskrit for self-control, then the Sanskrit for smoking, then say \emph{parityajati} again
\end{itemize}

\begin{tcolorbox}[breakable,colback=teal!4,colframe=teal!35!black,title={Before you move on}]
What does \sk{परित्यजति} mean? (\textquotedblleft{}Gives up entirely\textquotedblright{}.) Say it once more.
\end{tcolorbox}
\section[\texorpdfstring{\sk{मद्यम्} (madyam)}{मद्यम् (madyam)}]{\sk{मद्यम्} (madyam) — liquor}

\label{lesson:SA-C489-madyam}



\begin{quote}
Before the new one: say the Sanskrit for smoking, then the Sanskrit for gives up entirely.
\end{quote}

\subsection*{You'll want to know: \sk{मद्यम्}}
\textbf{\sk{मद्यम्}} — \emph{madyam} — \textquotedblleft{}liquor\textquotedblright{}.

\begin{quote}
\sk{मद्यं} \sk{यकृत्} \sk{पीडयति।}
\end{quote}

\begin{quote}
\emph{madyaṃ yakṛt pīḍayati.}
\end{quote}

\begin{quote}
Liquor harms the liver.
\end{quote}

A neuter noun.

\begin{cousinweb}
\textbf{\sk{मद्यम्}} (\emph{madyam}) comes from the root √mad, \textquotedblleft{}to be drunk, to rejoice\textquotedblright{}; MW gives \textquotedblleft{}any intoxicating drink, liquor, wine\textquotedblright{} (MW).
\end{cousinweb}

\begin{grammarlens}[title={What you've built}]
One more word for describing what you lived and felt.
\end{grammarlens}

\subsection*{Your turn}
\begin{itemize}
\raggedright
  \item \emph{Say it:} \emph{madyam}
  \item \emph{Say it:} \emph{madyam}, once more
  \item \emph{Say it:} \emph{madyaṃ yakṛt pīḍayati.} — Liquor harms the liver.
  \item \emph{Recall:} say the Sanskrit for smoking, then the Sanskrit for gives up entirely, then say \emph{madyam} again
\end{itemize}

\begin{tcolorbox}[breakable,colback=teal!4,colframe=teal!35!black,title={Before you move on}]
What does \sk{मद्यम्} mean? (\textquotedblleft{}Liquor\textquotedblright{}.) Say it once more.
\end{tcolorbox}
